\section{Limitations, Ethics, and Conclusion}

This study concerns one small public event stream. Windows overlap, policies change, official status
records can omit local failures, and community application reports are self-selected. The discovery
layer also depends on an external service, although final labels require primary evidence. Historical
point estimates should therefore not be generalized to other organizations or interpreted causally.

The design avoids claims about private motives, personality, or mental state. It uses public
professional actions, discourages attempts to manipulate the observed signals, and records possible
reflexive intervention. Rarity also depends on scale: the six-hour event rate is about 3.8\%, while
the daily rate is about 15\%. The central methodological result is modest: an operational-action forecast can be
made auditable by separating discovery, labels, features, locked predictions, outcomes, and
revisions. A recent-rate baseline outperforms the selected historical calendar model, and the tested
incident features add no stable value. Only the amended prospective sequence can determine whether
those orderings persist.
